\documentclass[10pt,a4paper]{amsart}
\usepackage[margin=19mm]{geometry}
\usepackage{amsmath,amssymb,mathtools}
\usepackage[scr=rsfs,cal=euler]{mathalfa}
\usepackage{xfrac}
\usepackage[unicode,hidelinks,bookmarks=false]{hyperref}
\newcommand{\ie}{i.e.\ }
\newcommand{\cf}{cf.\ }
\newcommand{\resp}{respectively\ }
\newcommand{\Subsection}{\subsection}
\providecommand{\nobreakdash}{\nobreak\hbox{-}}
\setlength{\emergencystretch}{2em}
\multlinegap0pt
\usepackage{bm}
\usepackage{bbm}
\usepackage{dsfont}
\usepackage{xspace}
\usepackage{cancel}
\usepackage{mathtools}
\usepackage{amsthm}
\makeatletter
\@ifundefined{definition}{\newtheorem{definition}{Definition}}{}
\@ifundefined{thm}{\newtheorem{thm}[definition]{Theorem}}{}
\makeatother
\title[Measures and dynamics on Pascal-Bratteli diagrams]{Measures and dynamics on Pascal-Bratteli diagrams}
\author{Sergey Bezuglyi, Artem Dudko, Olena Karpel}
\date{Source: arXiv:2411.06280v3 (CC BY 4.0)}
\begin{document}
\maketitle

\noindent\emph{Source and licence.} Measures and dynamics on Pascal-Bratteli diagrams. Version arXiv:2411.06280v3 (CC BY 4.0), distributed under the Creative Commons Attribution 4.0 International licence, as recorded in the arXiv licence metadata for this exact version and shown on the abstract page https://arxiv.org/abs/2411.06280v3. Full rights notice: licenses/arxiv-2411.06280-CC-BY-4.0.txt.\par\smallskip

\noindent\textit{Editorial scope.} Counted unit: the Thm whose source label is \texttt{thm: measure supports} in the article (source0.tex lines 420--427, source label \texttt{thm: measure supports}), delivered together with the article's own complete original proof of it (source0.tex lines 428--454, one \texttt{proof} environment of 27 lines). No mathematics is written, repaired or re-derived: statement and proof are the article's own text line for line. Required context retained uncounted: none. Every definition, display and label of those lines is retained unchanged. Omitted from this package and not redistributed: 1--419, 455--760. Nothing inside a retained excerpt is omitted or altered: every retained body is delivered exactly as the article's lines are written, so no omission receipt of any kind accompanies this package. Citations: the retained excerpts carry no citation command at all, so no bibliography entry of the article is transcribed and no reference is redistributed. Reference adaptation: none was needed and none was made; every reference inside the delivered unit resolves inside it, so no unresolved reference or citation is produced. Printed slips kept literal: no printed word, symbol, hypothesis or number of the counted statement or of its proof was altered. Layout and class: the article's own class is \texttt{amsart} with packages including inputenc, color, amsthm, amssymb, latexsym, amsmath, fontenc, MnSymbol, bbm, centernot, mathrsfs, amsfonts, textcomp, newlfont; the wrapper is the lane's pinned \texttt{amsart} editorial wrapper at 10pt on A4, which restates the theorem-like environments the retained text needs and adds the article's own macro definitions that the retained text needs (none), with extra wrapper packages (bm, bbm, dsfont, xspace, cancel, mathtools). The delivered numbering is the unit's own. Assets: no retained span contains a figure environment, a graphics-inclusion command, a PDF-page inclusion, a table or a TikZ picture; no image file is copied into this package, the package declares no asset dependency and no image file is redistributed with it. Licence: version arXiv:2411.06280v3 of this article is distributed under the Creative Commons Attribution 4.0 International licence, as recorded in the arXiv licence metadata for this exact version and shown on the abstract page https://arxiv.org/abs/2411.06280v3. A full rights notice is delivered at licenses/arxiv-2411.06280-CC-BY-4.0.txt. Characters: every retained excerpt is pure ASCII, so the delivered engine, XeLaTeX with its default Latin Modern text font, reports no missing character.
\begin{thm}\label{thm: measure supports}
For each $0 < p < 1$ and each $\varepsilon > 0$ there exists a subdiagram $B_p$  of the Pascal diagram such that 
\begin{enumerate} \item[$1)$] 
%\geqslant 
$\nu_p(X_p) > 1 - \varepsilon$ for each $p$, where $X_p$ is the path space of $B_p$; 
\item[$2)$] the subspaces $X_p,0\leqslant p\leqslant 1$, are pairwise disjoint.
\end{enumerate}
\end{thm}

\begin{proof} 

Fix $0<p<1$. We will construct by induction a sequence $N_i=N_i(p)$ such that for the sets 
$$A_i=A_i(p)=\left\{\{(x_n,y_n)\}\in X_B:\left|\frac{x_k}{k} -p\right|<2^{1-i}\;\;\text{for all}\;\;k > N_i\right\}$$ 
one has  $\nu_p(\bigcap_{i\leqslant j}A_i)>1 - \varepsilon$ for every $j\in\mathbb N$. Let $N_0=0$.
\vskip 0.1cm
\noindent {\bf Base of induction.} Since $0\leqslant x_n\leqslant n$ for all $n$ we set $A_0=X_B$. One has $\nu_p(A_0)=1.$
\vskip 0.1cm
\noindent {\bf Step of induction.} Let $j\in\mathbb N$. Assume that $N_i$ are constructed for $0\leqslant i\leqslant j$ such that 
$$
\nu_p\left(\bigcap_{i\leqslant j}A_i\right) > 1 - \varepsilon.
$$ 
By the Law of Large Numbers, for $\nu_p$ almost all $\{(x_n,y_n)\}\in X_B$ one has 
$$
\lim_{n \rightarrow \infty} \frac{x_n}{n} = p.
$$ 
It follows that for $N_{j+1}$ large enough one can make $\nu_p(A_{j+1})$ be arbitrary close to $1$. This implies that for $N_{j+1}$ large enough we have
$$\nu_p\left(\bigcap_{i\leqslant j+1}A_i\right)
%= \nu_p(A_{j+1}) - \nu_p\left(A_{j+1} \setminus \bigcap_{i\leqslant j}A_i\right) 
> 1 - \varepsilon,
$$
and that proves the induction step.

Now, let $A(p)=\bigcap_{i\in\mathbb Z_+}A_i(p)$. Introduce a subdiagram $B_p=(V_p,E_p)$ of the Pascal diagram as follows. Let $V_p$ be the set of all vertices $v\in \mathbb Z_+\times\mathbb Z_+$ such that there exists a path in $A(p)$ containing $v$. Let $E_p$ be the set of all edges of the Pascal diagram having both ends in $V_p$. 

Condition $1)$ of Theorem \ref{thm: measure supports} is satisfied by construction. Let $0 < p\neq q < 1$. Then there exists $i$ such that $|p-q|>2^{2-i}$. By definition, $A_i(p)$ and $A_i(q)$ are disjoint. This implies that condition $2)$ is satisfied as well.
\end{proof}

\end{document}
